\chapter{Exact phase-layer energies and geometric limits}\label{sec:examples-phase-geometry}
These examples refine the smooth-interface recovery calculation from \cite{ChAC}, Section 3.1, and the periodic ODE examples from Section \ref{sec:examples-phase-periodic}. They distinguish a particular competitor's energy from the existence or minimality of a critical point. The potential, $q$, $\psi$ and $c_W$ retain their normalization from Definition \ref{def:allen-cahn-energy}.

\begin{example}[Vertical graph layers retain an area excess]\label{ex:phase-vertical-graph-energy}
Let $\Omega\subset\mathbb R^d$ be bounded and open, and $h\in C^2(\Omega)\cap W^{1,\infty}(\Omega)$. On $\Omega\times\mathbb R$ put
$u_\varepsilon(y,z)=q((z-h(y))/\varepsilon)$. On this unbounded cylinder, energy and discrepancy mean the integrals of the local densities; the function is locally Sobolev but need not be globally in $L^2$. Then, for every $\varepsilon>0$,
\[
 E_\varepsilon(u_\varepsilon;\Omega\times\mathbb R)
  =c_W\int_\Omega\left(1+\frac{|Dh|^2}{2}\right)dy,
 \qquad
 \xi_\varepsilon^{u_\varepsilon}(\Omega\times\mathbb R)
  =\frac{c_W}{2}\int_\Omega|Dh|^2dy.
\]
The excess over $c_W$ times the graph area is exactly
\[
 \frac{c_W}{2}\int_\Omega
       \left(\sqrt{1+|Dh|^2}-1\right)^2dy.
\]
In particular this unnormalized vertical profile is generally not a recovery sequence for the graph perimeter.
\end{example}
\begin{proof}
Write $t=(z-h(y))/\varepsilon$. Then
$|Du_\varepsilon|^2=\varepsilon^{-2}\psi(t)^2(1+|Dh|^2)$ and $W(q(t))=\psi(t)^2/2$. The energy density is therefore
$\varepsilon^{-1}\psi(t)^2(1+|Dh|^2/2)$ and the discrepancy density is
$\varepsilon^{-1}\psi(t)^2|Dh|^2/2$. Integrate in $z$, using $dz=\varepsilon\,dt$ and $\int\psi^2=c_W$, to obtain both identities. The graph area is $\int_\Omega\sqrt{1+|Dh|^2}$ by the area formula. The elementary identity
$1+s/2-\sqrt{1+s}=(\sqrt{1+s}-1)^2/2$ proves the excess formula. It vanishes precisely when $Dh=0$ almost everywhere; the error does not disappear merely by taking a thinner transition layer.
\end{proof}

\begin{example}[The residual detects the same normalization error]\label{ex:phase-vertical-graph-residual}
For the preceding layer,
\[
 \varepsilon^2\Delta u_\varepsilon-W'(u_\varepsilon)
   =q''(t)|Dh|^2-\varepsilon\psi(t)\Delta h.
\]
If this is identically zero for all $(y,z)$, then $Dh=0$, so $h$ is constant on every connected component. When $h(y)=b\cdot y+c$, the correctly normalized tilted-plane solution is instead
\[
 \widetilde u_\varepsilon(y,z)
  =q\left(\frac{z-b\cdot y-c}{\varepsilon\sqrt{1+|b|^2}}\right).
\]
Its energy per base area is $c_W\sqrt{1+|b|^2}$ and its discrepancy vanishes.
\end{example}
\begin{proof}
Twice differentiating the composition gives
$\varepsilon^2\Delta u=q''(t)(1+|Dh|^2)-\varepsilon\psi(t)\Delta h$.
Subtract $W'(q)=q''$ to obtain the residual. Divide by $\psi>0$ at a fixed $y$: the factor $q''(t)/\psi(t)=-\sqrt2\tanh(t/\sqrt2)$ varies with $z$, whereas $Dh$ and $\Delta h$ do not. The identity for all $z$ forces $|Dh|^2=0$. For an affine graph the normalized argument is signed Euclidean distance divided by $\varepsilon$, so the planar-interface calculation of Proposition \ref{prop:phase-flat-interface} applies. Direct integration in $z$ gives the stated energy factor $\sqrt{1+|b|^2}$. The unit norm of the distance gradient gives equipartition pointwise.
\end{proof}

\begin{proposition}[The second-order energy of a distance layer in three-space]\label{prop:phase-tube-energy-correction}
Let $\Sigma\subset\mathbb R^3$ be a smooth compact embedded two-sided surface without boundary, and choose an injective symmetric tube $F(y,z)=y+z\nu(y)$, $|z|<\rho$, as in Lemma \ref{lem:phase-exact-fermi-geometry}. For $u_\varepsilon(F(y,z))=q(z/\varepsilon)$ define
\[
 m_2=\int_{\mathbb R}t^2\psi(t)^2dt\in(0,\infty).
\]
The energy on this tube satisfies
\[
 E_\varepsilon(u_\varepsilon;\operatorname{Tube}_\rho(\Sigma))
 =c_W\operatorname{Area}(\Sigma)
   +\varepsilon^2m_2\int_\Sigma K_\Sigma\,d\mu
   +O(e^{-2\sqrt2\rho/\varepsilon}).
\]
There is no linear term in $\varepsilon$.
\end{proposition}
\begin{proof}
The normal coordinate has unit gradient. Hence the energy density times the volume ratio is
$\varepsilon^{-1}\psi(z/\varepsilon)^2(1+Hz+K_\Sigma z^2)$.
Change variables $z=\varepsilon t$. On the symmetric interval the term $t\psi(t)^2$ is odd, so its integral is exactly zero. The remaining terms give the exact finite-interval formula
\[
 \operatorname{Area}(\Sigma)\int_{-\rho/\varepsilon}^{\rho/\varepsilon}\psi^2dt
 +\varepsilon^2\left(\int_\Sigma K_\Sigma\,d\mu\right)
     \int_{-\rho/\varepsilon}^{\rho/\varepsilon}t^2\psi^2dt.
\]
The exponential decay $\psi^2\leq Ce^{-2\sqrt2|t|}$ makes $m_2$ positive and finite. The omitted zeroth moment is $O(e^{-2\sqrt2\rho/\varepsilon})$. The second-moment tail, multiplied by $\varepsilon^2$, is bounded by $C(\rho^2+\varepsilon\rho+\varepsilon^2)e^{-2\sqrt2\rho/\varepsilon}$ and has the same stated order for fixed $\rho$. This proves the formula. It is an expansion for the specified layer on the tube; extending the function to an ambient domain requires a separate tail estimate such as the one in Theorem \ref{thm:phase-smooth-recovery}.
\end{proof}

\begin{example}[Spherical and toroidal corrections]\label{ex:phase-sphere-torus-corrections}
For a round sphere of radius $R$, the second-order term in Proposition \ref{prop:phase-tube-energy-correction} is $4\pi m_2\varepsilon^2$, independent of $R$. For every circular ring torus it is zero. This correction is governed by total Gauss curvature, whereas the Willmore energy of the same ring tori is positive and varies with the radius ratio.
\end{example}
\begin{proof}
On the sphere, $K_\Sigma=R^{-2}$ and its area is $4\pi R^2$, so $\int K_\Sigma=4\pi$. For a ring torus the already computed identity in Example \ref{ex:ring-torus-total-curvature} gives $\int K_\Sigma=0$. Insert these values in the tube formula. The ring-torus energy in Theorem \ref{thm:circular-ring-willmore-energy} is $\pi^2R^2/(r\sqrt{R^2-r^2})$. Thus no identification of the tube correction with the Willmore energy is justified by this example; the two curvature integrals are different.
\end{proof}

\begin{example}[An exact periodic solution converges to two interfaces]\label{ex:phase-circle-two-interface-limit}
Fix a circle of length $\ell$. For $0<\varepsilon<\ell/(2\pi)$ choose $a_\varepsilon$ by $T(a_\varepsilon)=\ell/\varepsilon$ and set
$u_\varepsilon(x)=f_{a_\varepsilon}(x/\varepsilon)$, normalized by the increasing zero at $x=0$. These are exact critical points, and
\[
 u_\varepsilon\longrightarrow
 \begin{cases}
 +1,&0<x<\ell/2,\\
 -1,&\ell/2<x<\ell
 \end{cases}
 \quad\text{in every finite }L^p(S^1_\ell).
\]
Their energy measures converge weakly to
$c_W(\delta_0+\delta_{\ell/2})$, and their total energy tends to $2c_W$.
\end{example}
\begin{proof}
The period equation implies $a_\varepsilon\to1$. For any fixed $b\in(0,1)$, the time for $f_a$ to rise from zero to $b$ converges, as $a\to1$, to
$\int_0^b(2W(s))^{-1/2}ds<\infty$. On the positive half-period the profile stays above $b$ except within this uniformly bounded time of either zero, by reflection about its maximum. Multiplication by $\varepsilon$ shows that $u_\varepsilon\to1$ uniformly on compact subintervals of $(0,\ell/2)$. The symmetry $f_a(t+T(a)/2)=-f_a(t)$ handles the negative half. Boundedness by one then gives every finite $L^p$ convergence.

Change variables $x=\varepsilon t$ in the energy: the total is exactly $\mathcal E(a_\varepsilon)\to2c_W$ by Proposition \ref{prop:phase-periodic-cell-energy}. Around either zero, for every fixed $L$ the rescaled profiles converge smoothly to a heteroclinic on $[-L,L]$. Their energy on the physical interval of radius $\varepsilon L$ therefore tends to $\int_{-L}^L\psi^2$. Sending $L\to\infty$ shows that every fixed neighborhood of either zero carries limiting energy at least $c_W$. Since the total limit is exactly their sum, energy outside any disjoint pair of such neighborhoods tends to zero. For a continuous test function, replace it on each small neighborhood by its value at the center and use uniform continuity. The resulting error tends to zero as the neighborhoods shrink, proving the stated weak convergence. This limit has two perimeter points of multiplicity one, even though each approximating critical point is unstable.
\end{proof}
